% Paul E. West and Sean T. Hayes

\documentclass[14pt,aspectratio=1610]{beamer}
% \documentclass[14pt,aspectratio=1610, handout]{beamer}
\usepackage{csu-theme}

\title{\Huge{}Testing in {\ttfamily\bfseries C++}}
\subtitle{CSCI 315: \textit{Data Structure Analysis}}
\institute{Charleston Southern University}

%% Comment out date to use default (today)
% \date{January 23, 2025}
\subject{Software Programming}
%\keywords{}

\begin{document}

\begin{frame}
  \titlepage{}
\end{frame}

\section*{Ethics Essay}

\begin{frame}{Ethics Essay on the Importance of Testing}
\centering\Large
\href{https://github.com/csu-cs/CSCI-315-2026-Fall/blob/master/ethics/README.md}{Let's go over your ethics essay today!}
\end{frame}

\section{Test Your Code}

\begin{frame}{Learning Outcomes: Testing}
  Students will be able to:
  \begin{overlayarea}{\linewidth}{9\baselineskip}
    \begin{enumerate}
      \item \emph{Explain}: Articulate the purpose of testing, differences between unit testing and other testing types, and why tests improve developer productivity.
      \item \emph{Design Tests}: Design effective unit tests including normal cases, edge cases, and failure cases.
      \item \emph{Write Tests}: Write unit tests using Doctest and use common test macros/assertions.
    \end{enumerate}
  \end{overlayarea}
\end{frame}

\begin{frame}{Learning Outcomes: Testing}
  Students will be able to:
  \begin{overlayarea}{\linewidth}{9\baselineskip}

    \begin{enumerate}\addtocounter{enumi}{3}
      \item \emph{Interpret Results}: Run tests, read test output, and diagnose failing tests to locate defects.
      \item \emph{Automate}: Integrate tests into a simple automated workflow using Makefiles.
      \item \emph{Use Test-First Practices}: Apply basic test-driven development to guide incremental implementation.
      \item \emph{Debug \& Prevent Regressions}: Create regression tests for fixed bugs and use tests to prevent regressions.
    \end{enumerate}
\end{overlayarea}
\end{frame}

\begin{frame}{Testing}
\textbf{\large\textit{Testing will make you more productive.}}\\[1em]
Meaning, you will spend less time working on problems because you will\textellipsis{}
\begin{itemize}
\item
  increase your chances of getting it right the first time AND
\item
  rarely make the same mistake twice.
\end{itemize}
\end{frame}

\begin{frame}{Warning}
\begin{itemize}[<+->]
\item A computer program \textit{cannot} sufficiently generate \\test cases for your code.
\item \href{https://en.wikipedia.org/wiki/Halting_problem}{The Halting Problem}: Alan Turing proved that it is impossible for a computer to determine if any program halts (finishes).\\
(\href{https://youtu.be/92WHN-pAFCs?t=14}{see Explanation on YouTube.})
\item If we don't know if it halts, how can we test?
\item This is your job, and an advantage you have over computation.
\end{itemize}
\end{frame}

\begin{frame}{Unit Test}
\begin{itemize}
\item
  For this course, unit testing is sufficient.
\item
  There are many other forms of testing, but this\\ course focuses on (relatively) small units of code.
\item
  There are numerous unit-testing tools, but no \textit{de facto} standard for C++.
  \begin{itemize}
  \item Unlike Java's JUnit
  \end{itemize}
\item
  Popular testing frameworks for C++ include \href{https://google.github.io/googletest/}{GoogleTest}, \href{https://www.boost.org/doc/libs/1_81_0/libs/test/doc/html/index.html}{Boost.Test}, \href{CppUnit}{CppUnit}, \href{https://cxxtest.com/}{CxxTest}, and \href{https://github.com/catchorg/Catch2}{Catch2}.
\end{itemize}
\end{frame}

\begin{frame}{Doctest -- Unit Testing for C++}
\begin{itemize}
\item
  We will use \href{https://github.com/doctest/doctest\#readme}{Doctest} for this course.
  \begin{itemize}
  \item
    Doctest is simple, yet has a rich feature set.
  \end{itemize}
\item
  You only need to place the one file, \href{https://github.com/doctest/doctest/blob/master/doctest/doctest.h?raw=true}{doctest.h}, in your test repository.
\item
  I have provided a full template in class-code.\\Let's take a look\textellipsis{}
\end{itemize}
\end{frame}

\section*{Lab 5}

\begin{frame}{Auto-Grader}
  \begin{itemize}
  \item An auto-grader will provide a grade for most of\\ the labs (starting in Lab 5).
  \begin{itemize}
    \item It will not say what test cases you failed.
    \item It may not mention if there is a timeout.
    \item It may not say if there are compilation errors.
    \item It \textit{may} tell you which part of the lab you got right or wrong.
  \end{itemize}
  \item Test cases from the auto-grader will \textbf{NOT} be disclose.  You \textbf{must} learn how to write your own test cases!
  \item The grade from the auto-grader should be pushed to your repository shortly after you push your solution.
  \end{itemize}
\end{frame}

\begin{frame}{Lab 5: A Review of C++ Arrays}
\centering\Large
Let's talk about Lab 5.
\end{frame}

\end{document}
